\section{Schnittstellen \& Peripherie}

\subsection[uC-Peripherie]{$\mu$C-Peripherie}
\begin{itemize}
    \item Ankopplung über Systembus des $\mu$C's.
    \item Hardwaremässige Umsetzung.
    \item Control-, Daten- \& Statusregister.
    \item Ein- bzw. Ausgabe über I/O-Ports.
\end{itemize}

\subsection{Schieberegister}
\textbf{Parallel In Serial Out (PISO):}
Einfache Verkettung von Flip-Flops. Zusätzliche Multiplexer nötig.

\textbf{Serial In Parallel Out (SIPO):}
Einfache Verkettung von Flip-Flops.

\vspace{0.5em}
\noindent\framebox{\parbox{\linewidth}{\centering \vspace{10pt}\textit{[Placeholder: Diagramme zu PISO / SIPO / Peripherie Blockschaltbild]}\vspace{10pt}}}
\vspace{0.5em}


\subsection[Taktgeber in einem uP-System]{Taktgeber in einem $\mu$P-System}
Wichtigkeit, den Taktgeber richtig zu wählen:
\begin{itemize}
    \item Für Synchronität, Genauigkeit (Präzision), Strom/Effizienz.
    \item Clock soll so gewählt sein, dass alle zeitkritischen Anforderungen in der definierten Zeit abgearbeitet werden können.
    \item Performance vs. Leistungsverbrauch.
\end{itemize}

\subsection{Direct Memory Access (DMA)}
\begin{itemize}
    \item Daten werden nicht mittels Load/Store über CPU transferiert, sondern direkt über den DMA.
    \item Daten werden nicht bearbeitet oder verändert, sie werden als ganze Pakete transferiert.
    \item Erst am Ende bzw. nach der Hälfte eines Datenpaket-Transfers wird ein Interrupt generiert.
    \item Interrupt Stacking \& Fetching wird reduziert.
\end{itemize}

\textbf{Anwendungsbeispiele:}
\begin{itemize}
    \item \textit{High Performance:} Datentransfer \& Verarbeitung erfolgt in grossen Paketen. CPU kann Daten verarbeiten während der DMA transferiert.
    \item \textit{Ultra Low Power:} Datentransfer \& -Verarbeitung erfolgt in grossen Paketen. CPU ist im Standby Modus während der DMA transferiert.
\end{itemize}


\subsection{UART}
\begin{itemize}
    \item Besitzt keinen Clk. (asynchrone Übertragung).
    \item Ausschliesslich Punkt-zu-Punkt Verbindung (Tx, Rx).
    \item Full-Duplex Betrieb uneingeschränkt möglich.
    \item Diverse Treiber (RS-232, RS-422, RS-485) für längere Distanzen.
    \item Stromsparend (Push-Pull-Ausgänge).
\end{itemize}
\textbf{Synchronisation:} Mittels des Start-Bits, welches bei jedem Symbol zu Beginn übertragen wird.\\
\textbf{UART vs. RS-232:}
UART: Logik-Pegel $1=\text{High}$, $0=\text{Low}$. \\
RS-232: Schnittstellen-Pegel $1=-3V \dots -12V$, $0=+3V \dots +12V$.

\subsection{SPI}
Synchrone Schnittstelle mit minimalen Spezifikationen.
\begin{itemize}
    \item Durch Synchronität höhere Übertragungsraten als mit UART möglich.
    \item Eingeschränkter Full-Duplex Betrieb oder Half-Duplex.
    \item Leitungen: CS, SCLK, MOSI, MISO.
    \item Mehrere Slaves am gleichen Interface. Nur kurze Distanzen möglich.
    \item Stromsparend (Push-Pull-Ausgänge).
\end{itemize}

\textbf{SPI: SW Data-Ready bei Subs}
\begin{itemize}
    \item Problematik: Sub kann dem Main kein "Data Ready" senden, Main muss aktiv ein Status-Register abfragen.
    \item Vorteil: Keine zusätzlichen Leitungen.
    \item Nachteile: Ständige Abfrage des Status-Registers, erhöhter Datentransfer, kein effizienter Standby-Betrieb möglich.
\end{itemize}

\textbf{SPI: HW Data-Ready bei Subs}
HW-Abfrage: Zusätzliche Interrupt-Leitungen und Data-Ready-Leitung.
\begin{itemize}
    \item Vorteile: Keine überflüssiger Datentransfer, effizienter Standby-Betrieb möglich.
    \item Nachteil: Zusätzliche Leitungen zu SPI-Schnittstelle.
\end{itemize}

\subsection{I2C}
Synchrone Schnittstelle mit höheren Spezifikationen.
\begin{itemize}
    \item Half-Duplex Bus Topologie.
    \item Nicht stromsparend (Open-Drain-Ausgänge mit Widerstand).
\end{itemize}